% Clase 4 --- Adelanto: la raíz cuadrada con regla y compás (§8.3).
% J es el simétrico de LA UNIDAD I respecto de O (J representa -1; el
% docente lo dice explícito en la transcripción: "el simétrico... que
% corresponde a menos uno"). El segmento J-X (de longitud 1+x) es diámetro
% de una semicircunferencia; O divide ese diámetro en JO=1 y OX=x. La
% perpendicular a la recta en O corta la semicircunferencia en Z, con
% OZ=raíz(1*x)=raíz(x) por el teorema de la altura (media geométrica) en
% el triángulo rectángulo J-Z-X inscrito en la semicircunferencia.
% Valor ilustrativo: x=4 (para que raíz(x)=2 sea verificable a ojo).
\begin{center}
\begin{tikzpicture}[scale=1.05]

  % ---- recta principal ----
  \draw[figaxis] (-1.7,0) -- (4.7,0) node[right, figlbl] {$R$};

  % ---- puntos clave sobre la recta ----
  \node[figptocerrado] (J) at (-1,0) {};
  \node[figptocerrado] (O) at (0,0)  {};
  \node[figptocerrado] (I) at (1,0)  {};
  \node[figptocerrado] (X) at (4,0)  {};

  \node[figlbl, below=3pt] at (J) {$J$};
  \node[figlbl, below=3pt] at (O) {$O$};
  \node[figlbl, below=3pt] at (I) {$I$};
  \node[figlbl, below=3pt] at (X) {$X$};

  % ---- semicircunferencia de diámetro JX (centro M, radio 2.5) ----
  \draw[figline] (1.5,0) ++(180:2.5) arc (180:0:2.5);

  % ---- Z: intersección de la perpendicular en O con la semicircunferencia ----
  \coordinate (Z) at (0,2);
  \node[figptocerrado, draw=figred, fill=figred] at (Z) {};
  \node[figlbl, figred, above=3pt] at (Z) {$Z$};

  % ---- perpendicular O-Z, remarcada como el segmento resultado ----
  \draw[figred, line width=1.5pt] (O) -- (Z);
  \node[figlbl, figred, right=2pt] at (0,1) {$\sqrt{x}$};

  % ---- triángulo rectángulo inscrito J-Z-X (auxiliar, justifica el "por qué") ----
  \draw[figaux] (Z) -- (J);
  \draw[figaux] (Z) -- (X);

  % ---- marca de ángulo recto en Z (cuadrado pequeño, lados hacia J y X) ----
  \coordinate (ru) at ($(Z)!{0.18/(2*sqrt(5))}!(X)$);
  \coordinate (rv) at ($(Z)!{0.18/sqrt(5)}!(J)$);
  \coordinate (rc) at ($(ru)+(rv)-(Z)$);
  \draw[figgray] (rv) -- (rc) -- (ru);

\end{tikzpicture}\\[0.5em]
{\small\itshape Media geométrica con regla y compás: $J$ es el simétrico de
la unidad $I$ respecto de $O$ (representa $-1$), de modo que el segmento
$JX$ --- diámetro de la semicircunferencia --- mide $1+x$ y $O$ lo divide en
$JO=1$ y $OX=x$. Como $JZX$ es un triángulo rectángulo inscrito en la
semicircunferencia (ángulo recto en $Z$), el teorema de la altura da
$OZ=\sqrt{JO\cdot OX}=\sqrt{x}$ (aquí, a modo de ilustración, $OX=4$ y
$OZ=2$). Válido solo como anticipo geométrico: la existencia de $\sqrt{x}$
recién queda justificada con el axioma de completitud (próxima clase).}
\end{center}
